\documentclass[11pt,a4paper]{article}
\usepackage[margin=25mm]{geometry}
\usepackage{amsmath,amssymb,amsthm,hyperref}
\setlength{\emergencystretch}{2em}
\newtheorem{theorem}{Theorem}
\newtheorem{lemma}[theorem]{Lemma}
\newtheorem{proposition}[theorem]{Proposition}
\newcommand{\F}{\mathbf F}
\newcommand{\gen}[1]{\langle#1\rangle}
\newcommand{\Om}{\Omega}
\title{A free endpoint for odd orthogonal strings in dimension five}
\author{Human Compression project: F001 working component}
\date{}
\begin{document}\maketitle
This is a construction and exceptional-field component, not yet a complete recovery of
Brooksbank's rank spectrum theorem. We prove existence in odd fields
of sufficiently large order without the primitive-parameter density
assertion in the supplied preprint. Exact finite certificates supply
the remaining odd construction cases and the exceptional field
of order three. The upper bound on ranks for the other odd fields
is a separate obligation. This working edition has not received
distinct review.

\section{Classical inputs and the three-dimensional step}
Let $F=\F_q$ have odd order, let $Q$ be a quadratic form, and write
$B(u,v)=Q(u+v)-Q(u)-Q(v)$. For $Q(u)\ne0$ its reflection is
\[
 s_u(v)=v-\frac{B(v,u)}{Q(u)}u.
\]
For a subspace $U$ write $K_\square(U)$ for the group generated by
the reflections whose roots lie in $U$ and have nonzero square norm;
all such groups act on the same ambient space. We retain the
classical spinor-norm generation theorem: on a nondegenerate
space, these reflections generate the kernel of the spinor norm.
We also retain the three-dimensional identifications
$\mathrm{SO}_3(q)\cong\mathrm{PGL}_2(q)$ and
$\Om_3(q)\cong\mathrm{PSL}_2(q)$, and Dickson's subgroup
classification in the following exact form, as stated in
Leemans--Moerenhout--O'Reilly-Regueiro, \emph{Projective linear
groups as automorphism groups of chiral polytopes}, supplied
twenty-five-page edition, Proposition 2.3, p.5. A subgroup of
$\mathrm{PSL}_2(q)$ or $\mathrm{PGL}_2(q)$ is either contained
in a point stabilizer, a torus normalizer, a subfield subgroup,
one of the exceptional groups $A_4,S_4,A_5$, or contains
$\mathrm{PSL}_2(q)$. The full proofs of these classical inputs
remain external acquisition costs. The string end-parabolic
criterion is the already acquired McMullen--Schulte criterion.

Put $\epsilon=\chi(-1)$ and $N=(q+\epsilon)/2$, where $\chi$
is the quadratic character. Thus $N$ is odd. Call a nondegenerate
plane $L$ suitable if $K_\square(L)$ is dihedral of order $2N$.
Its rotation subgroup then has order $N$.

\begin{lemma}\label{three}
Suppose $q\ge13$. Let $T=L+\gen x$ be a three-space containing
a suitable plane $L$, where $Q(x)$ is a nonzero square and
$x\notin L\cup L^\perp$. Then
\[
                  \gen{K_\square(L),s_x}=K_\square(T).
\]
The conclusion holds whether or not $T$ is nondegenerate, provided
the ambient space is nondegenerate.
\end{lemma}
\begin{proof}
First suppose $T$ nondegenerate. Let $D=K_\square(L)$ and let
$C\le D$ be its cyclic subgroup of order $N$. Under
$g\mapsto\det(g)g$, the group $K_\square(T)$ maps onto either
$\mathrm{PSL}_2(q)$ or $\mathrm{PGL}_2(q)$, with kernel
$\{1,-1_T\}$ or $1$, respectively. This follows directly from
the determinant and spinor norm: the first case occurs precisely
when the spinor norm of $-1_T$ is square. The map is injective
on $D$, since $D$ fixes the nonzero perpendicular line $L^\perp$
pointwise. Hence its image contains a dihedral group of order $2N$.

Here $N>5$ and $N>\sqrt q+1$. The exceptional groups and proper
subfield groups cannot contain an element of order $N$: the
semisimple orders in $\mathrm{PGL}_2(q_0)$ are at most $q_0+1$,
where $q_0\le\sqrt q$. A point stabilizer cannot contain this
dihedral group either. Its image on its diagonal quotient is
abelian and injective on $C$, whereas a reflection of $D$ inverts
$C$. Dickson's classification leaves only the torus normalizer
and a group containing $\mathrm{PSL}_2(q)$.

The normalizer of $C$ in the three-dimensional orthogonal group
preserves $L$ and $L^\perp$: these are its nontrivial and trivial
constituents. A reflection preserves $L$ precisely when its root
lies in $L\cup L^\perp$. Thus $s_x$ excludes the normalizer.
If the target is $\mathrm{PGL}_2(q)$, the image of $s_x$ has
nonsquare spinor norm, so the image is the whole target rather
than its index-two subgroup. If the target is
$\mathrm{PSL}_2(q)$, the full group upstairs is
$\gen{-1_T}\times\mathrm{PSL}_2(q)$; a subgroup mapping
bijectively onto the second factor is a graph of a homomorphism
to $C_2$. Perfection of $\mathrm{PSL}_2(q)$ makes that graph
have determinant one, contrary to the presence of $s_x$.
In both cases the generated group is $K_\square(T)$.

Now suppose $T$ degenerate. Write $T=L\perp\gen z$ with
$Q(z)=0$, and choose a singular partner $w$ with $B(z,w)=1$,
perpendicular to $L$. The group supported on $T$ has translation
kernel $\{t_u:u\in L\}$, where
\[
 t_u(y)=y-B(y,u)z\ (y\in L),\quad
 t_u(z)=z,\quad t_u(w)=w+u-Q(u)z.
\]
These maps multiply by $t_ut_v=t_{u+v}$. The subgroup generated
by square-norm reflections is $L^+\rtimes D$: the displayed
translations have odd order and trivial spinor norm, and products
of the reflections in $u$ and $u+az$ provide all scalar multiples
of every square-norm $u$ as translation parameters. Such roots
span $L$; alternatively this follows from the irreducibility below.
Writing $x=u+az$, with $a\ne0$, shows that $s_xs_u$ is a nonzero
translation. Its $D$-conjugates supply the additive span of a
nonzero $D$-orbit.

That span is all of $L$, even over the prime field. Here are the
field details. Write $q=p^e$, and let $\lambda$ be an eigenvalue
of a generator of $C$, of order $N$. In the nonsplit case its
degree over $\F_p$ is $2e$: any proper degree divides $2e$ and
is at most $e$; degree $e$ would put $N$ into $q-1$, and a
smaller degree is impossible because
$N>p^{e-1}-1$. In the split case $\lambda$ has degree $e$.
The two eigenvalues are not conjugate over $\F_p$, since such
conjugacy would give $N\le p^{e-1}+1\le q/3+1<N$ for $q>9$.
Their two irreducible prime-field constituents are exchanged by
the reflections in $D$. Thus $D$ acts irreducibly on $L$ over
$\F_p$ in either case. The nonzero translation therefore supplies
the whole kernel and proves the assertion.
\end{proof}

\begin{lemma}[Saturation from a suitable plane]\label{sat}
Suppose $q\ge13$, $V$ is nondegenerate of dimension at least four,
and $W$ is a nondegenerate hyperplane. Let $L$ be a suitable plane
not contained in $W$, whose intersection with $W$ has square norm.
Then $\gen{K_\square(W),K_\square(L)}=K_\square(V)$.
\end{lemma}
\begin{proof}
It suffices to obtain $s_y$ for every square-norm root
$y\notin W\cup L$. Put $T=L+\gen y$ and $M=T\cap W$.
The two-space $M$ contains the square-norm line $L\cap W$.
If $M$ is nondegenerate, it has $(q-1)/2$ or $(q+1)/2$
square-norm projective points. If it is degenerate, the presence
of that line shows it has $q$ such points. There are therefore
at least six. At most one lies in $L$, and at most one in
$L^\perp\cap T$. Choose another, represented by $x$.
Its reflection is already in $K_\square(W)$; Lemma~\ref{three}
gives $K_\square(T)$ and in particular $s_y$.
\end{proof}

\section{A path with one free bond}
Consider the five-space with quadratic form
\[
 Q(X)=\sum_{i=1}^5 X_i^2+
       \eta X_1X_2+t(X_2X_3+X_3X_4+X_4X_5),       \tag{1}
\]
where $t\eta\ne0$. Suppose every consecutive coordinate subspace
is nondegenerate, and each of the final three adjacent products
of coordinate reflections has order $N=(q+\chi(-1))/2$.
The corresponding two-space groups are then precisely their
$K_\square$ groups: two reflections of square norm generate a
dihedral group of order $2N$, the full spinor kernel in that plane.

\begin{proposition}\label{path}
For $q\ge13$, the five coordinate reflections of (1) form a
string C-group and generate $K_\square(V)$. If the determinant
of the symmetric matrix representing $Q$ is a nonzero square,
their negatives form a rank-five string C-group for $\Om_5(q)$.
It has rank-four and rank-three contractions.
\end{proposition}
\begin{proof}
Every consecutive triple has a suitable pair. Its other root is
outside that plane and is not perpendicular to it, so
Lemma~\ref{three} gives the full $K_\square$ group on the triple.
For an interval $v_1,\ldots,v_b$ of length $b\ge4$, its suffix
$W=\gen{v_2,\ldots,v_b}$ has full $K_\square(W)$ by induction.
The plane $L_0=\gen{v_2,v_3}$ is suitable. Conjugate it by
$s_{v_1}$ to obtain a suitable plane $L$ not contained in $W$.
The line $L\cap W=\gen{v_3}$ has norm one. Lemma~\ref{sat}
therefore gives the full group on the interval. This proves
generation for all intervals of length at least three.

Every two-generator subgroup is a dihedral C-group. For the
intersection step, an isometry supported on a nondegenerate
subspace $U$ fixes $U^\perp$ pointwise; conversely, its restriction
to $U$ determines it. Thus for nondegenerate $U,W,U\cap W$,
\[
 K_\square(U)\cap K_\square(W)=K_\square(U\cap W).
\]
The same assertion on a square-norm one-space is its order-two
reflection group. For a triple this bounds the intersection of
the two dihedral end groups by its middle reflection. For length
four the common middle plane is suitable, and for length five
the common middle three-space has already been generated.
Induction with the end-parabolic criterion proves the full
intersection property. Nonadjacent coordinate roots are
perpendicular, giving the string relations.

The spinor norm of $-1_V$ is the determinant of $Q$ modulo
squares. Under the determinant hypothesis it lies in
$K_\square(V)$, and $g\mapsto\det(g)g$ maps this group onto
$\Om_5(q)$ with kernel $\gen{-1_V}$. Each end parabolic fixes
a nonzero perpendicular line pointwise, so the map is injective
there. We must also exclude a new intersection after quotienting.
Let $a$ and $b$ span the perpendicular lines to the two end
hyperplanes. Their pairing is nonzero: the corner entry of the
inverse tridiagonal form matrix is its four nonzero adjacent
entries multiplied together, divided by its nonzero determinant.
If elements $h_0,h_4$ of the two end groups had $h_0=-h_4$,
then $h_0$ would fix $a$ and negate $b$. This contradicts
$B(a,b)\ne0$ in odd characteristic. Thus the intersection of
the two image groups is exactly the image of their old intersection.
Their individual markings are faithful images of C-groups, so
the end-parabolic criterion proves the assertion. Their
adjacent product orders equal those of the original reflections.
The first adjacent pair is noncommuting because its independent
roots have nonzero pairing $\eta$; the remaining pairs have order
$N>2$. The marking is therefore irreducible. The third and fourth
bonds have odd order $N$, so two applications
of the acquired rank-contraction theorem give ranks four and three.
\end{proof}

\section{Uniform choice of parameters}
\begin{theorem}\label{uniform}
If $q\ge33$ is an odd prime power, then $\Om_5(q)\cong
\mathrm{PSp}_4(q)$ has string C-group representations of ranks
five, four and three.
\end{theorem}
\begin{proof}
Let $\epsilon=\chi(-1)$, $N=(q+\epsilon)/2$, and take an element
$\mu$ of exact order $N$ in $\F_q^\times$ when $\epsilon=-1$,
or in the norm-one torus of $\F_{q^2}^\times$ when
$\epsilon=1$. Put $t=\mu+\mu^{-1}\in F$ and $x=t^2/4$.
The product of reflections on a plane with norm-one roots and
pairing $t$ has characteristic polynomial
$X^2-(t^2-2)X+1$ and eigenvalues $\mu^2,\mu^{-2}$.
It therefore has order $N$. Also
\[
 1-x=-\frac{(\mu-\mu^{-1})^2}{4}
\]
is a nonsquare in $F$ in both torus cases. In particular $t$ and
$1-x$ are nonzero. The equality $1-2x=0$ would imply
$\mu^4=-1$, impossible for the odd-order element $\mu$.

Choose $\mu$ so that also $A=1-3x+x^2\ne0$.
The equation $A=0$, multiplied by $16\mu^4$, is
\[
 \mu^8-8\mu^6-2\mu^4-8\mu^2+1=0.                 \tag{2}
\]
There are at most eight roots, while $\varphi(N)>8$ since
$N\ge17$ is odd. For completeness, the odd integers with
Euler totient at most eight are exactly $1,3,5,7,9,15$:
their prime divisors can only be $3,5,7$, and multiplicativity
and the prime-power formula give this list. Hence a permissible
primitive $\mu$ exists without any density estimate.

Now let $\eta$ vary in $F$. The determinants of the three
proper prefixes of length two, three and four are, with
$y=\eta^2/4$,
\[
 1-y,\qquad 1-x-y,\qquad 1-2x-y(1-x).             \tag{3}
\]
The suffix determinants of lengths two, three and four are
$1-x$, $1-2x$, and $A$, all already nonzero. The full determinant
is
\[
 \Delta=A-B\eta^2,\qquad B=(1-2x)/4\ne0.          \tag{4}
\]
For $A,B\ne0$, the elementary character identity
\[
 \sum_{\eta\in F}\chi(A-B\eta^2)=-\chi(-B)
\]
follows by counting $(u-v)(u+v)=c$ for $c\ne0$.
At most two summands are zero. Consequently at least
$(q-3)/2$ values of $\eta$ give a nonzero square $\Delta$.
Discard $\eta=0$ and the at most six roots in (3). Since
$(q-3)/2>7$, a value remains. Every consecutive subspace is
then nondegenerate, and Proposition~\ref{path} applies.
\end{proof}

\section{Finite parameters and the four small fields}
The hypothesis $q\ge33$ in Theorem~\ref{uniform} is used only
to choose the parameters. Proposition~\ref{path} applies already
for $q\ge13$. The following table supplies the intervening fields.
Here $z_{25}$ is a root of $X^2-X+2$ over $\F_5$, and $z_{27}$
is a root of $X^3-X+1$ over $\F_3$, chosen as in the accompanying
GAP field model. Equivalently these are its Conway generators
$Z(25)$ and $Z(27)$; the certificate states the exact powers.
\[
\begin{array}{c|rrrrrrrr}
 q&13&17&19&23&25&27&29&31\\\hline
 t&10&14&3&14&z_{25}^{11}&z_{27}^{16}&14&16\\
 \eta&1&3&4&1&z_{25}^{2}&z_{27}^{4}&4&19
\end{array}                                                    \tag{5}
\]
The finite checker verifies all fifteen interval determinants,
the square class of the full determinant, and the exact order
$(q+\chi(-1))/2$ of the last three products for every row.
Thus no group enumeration is needed for these eight fields.

For $q=5,7,9,11$ use
\[
 Q(X)=\sum_{i=1}^5X_i^2+\sum_{i=1}^4c_iX_iX_{i+1}
\]
with the following coefficients. The element $z\in\F_9$ satisfies
$z^2=z+1$ and has order eight.
\[
\begin{array}{c|c|c|r}
q&(c_1,c_2,c_3,c_4)&\text{adjacent orders}&|\Om_5(q)|\\\hline
5 &(1,1,2,1)&(3,3,5,3)&4\,680\,000\\
7 &(1,1,1,2)&(3,3,3,7)&138\,297\,600\\
9 &(z,1,1,z)&(5,3,3,5)&1\,721\,606\,400\\
11&(1,1,1,3)&(3,3,3,5)&12\,860\,654\,400
\end{array}                                                    \tag{6}
\]
In these four cases some proper intervals may be degenerate, so
we use the complete finite certificate in place of
Proposition~\ref{path}. Let $M$ be the symmetric matrix of $Q$.
The reflection in the $i$th coordinate root is the row-action
matrix $R_i=I-2(Me_i)e_i^{\mathsf T}$; its negative has
determinant one. The certificate verifies $R_iMR_i^{\mathsf T}=M$
and that $\det M$ is a nonzero square. Hence the negative
reflections lie in $\Om_5(q)$ by their spinor norms.

Their permutation action is computed on every isotropic
projective point, normalized by making the first nonzero
coordinate one. There are $(q+1)(q^2+1)$ such points, namely
156, 400, 820 and 1464 in the four cases. This action is faithful
on the group of negative reflections. The certificate exhibits
five independent isotropic representatives and verifies that
the supports of the remaining representatives connect all five
coordinates in that basis. A linear map fixing all their
projective points is diagonal in the chosen basis, and those
connected supports make all five diagonal entries equal. A
scalar isometry of determinant one in odd dimension is the
identity. This proves the required faithfulness without a
faithfulness assumption about a permutation library model.

The permutation groups have the orders in (6), equal to
$q^4(q^2-1)(q^4-1)/2$, the classical order of $\Om_5(q)$.
They are therefore the full groups. Each tuple has five
involutions, the displayed adjacent orders, and the string
commutations. For all $32^2=1024$ ordered subset pairs the
certificate verifies the intersection condition. Containment
pairs are immediate; other pairs are checked by exact group
intersection, with the symmetric pair reused. The last two
bonds in every row have odd order, so the acquired contraction
theorem gives ranks four and three as well.

\begin{theorem}\label{positive}
For every odd prime power $q\ge5$, the group
$\mathrm{PSp}_4(q)\cong\Om_5(q)$ has string C-group
representations of ranks three, four and five.
\end{theorem}
\begin{proof}
Combine Theorem~\ref{uniform}, the checked rows (5), and the
four exact constructions (6).
\end{proof}

\section{The exceptional group of order 25920}
We give a separate exhaustive certificate that
$\mathrm{PSp}_4(3)$ has no string C-group representation.
To identify the group without naming a library group, use the
alternating form
\[
 J=\begin{pmatrix}0&1&0&0\\-1&0&0&0\\
                  0&0&0&1\\0&0&-1&0\end{pmatrix}
\]
on $\F_3^4$ and the transvections
$w\mapsto w+(wJv^{\mathsf T})v$ for the forty normalized
projective vectors $v$. Their action on these same forty
points has order 25920, half the classical order of
$\mathrm{Sp}_4(3)$. Thus it is the full projective symplectic
group. It is simple, as also verified by the finite checker,
and has 315 involutions in two classes of sizes 45 and 270.

Here is the complete extension procedure. A node is an ordered
prefix $(s_1,\ldots,s_r)$ that is already a connected string
C-marking of its generated subgroup $H$. Fix one representative
for $s_1$ in each involution class. At a node take the involutions
\[
 x\in C_G(s_1,\ldots,s_{r-1})\setminus H,
       \qquad [x,s_r]\ne1,                              \tag{7}
\]
modulo conjugation by $C_G(s_1,\ldots,s_r)$. For each orbit
representative append $x$ and test the intersection property
by the recursive end-parabolic criterion. Keep precisely those
extensions that pass. The centralizers in (7) and their orbits
are calculated in the explicit forty-point group.

This procedure is exhaustive. A string C-marking is irredundant,
so its next generator cannot lie in $H$. The string relations
are exactly the centralizer condition in (7), and conjugating by
the pointwise stabilizer of the prefix preserves every extension
property. A marking for a nonabelian simple group must be connected:
a missing bond would split its generators into nontrivial
commuting normal subgroups. Its every prefix is itself a C-marking,
so discarding a failed prefix cannot remove a successful full
marking. Finally every retained extension at least doubles $|H|$.
No irredundant sequence in a group of order 25920 has length more
than fourteen, giving a termination bound independent of the
observed early termination.

The complete counts are
\[
\begin{array}{c|rrrrrr}
\text{length}&1&2&3&4&5&6\text{--}14\\\hline
\text{retained prefixes}&2&13&121&20&2&0\\
\text{rejected by intersection}&0&0&20&25&0&0
\end{array}                                                     \tag{8}
\]
None of the retained prefixes generates $G$. Therefore there is
no connected string C-marking of $G$ at any rank, and simplicity
excludes the disconnected ones.

\paragraph{Attribution and exact remaining scope.}
The symmetry construction, spinor-norm passage and desired rank
spectrum come from Brooksbank, \emph{On the ranks of string C-group
representations for symplectic and orthogonal groups},
arXiv:1912.05512v1, 11 December 2019, Theorem 1.1 and Sections 2--4.
The endpoint method reuses our acquired characteristic-two
endpoint proof for L031; the prime-field translation argument,
odd-characteristic generation step and the parameter count above
are the present contribution. The rank-contraction theorem is the
acquired L037 result. The nonorthogonal perpendicular lines give
the separate intersection check required by the central quotient;
endpoint injectivity alone is not used to infer it. We retain,
rather than reprove, the
classical facts listed in the first section and the orders of the
classical groups used to identify the finite examples. The finite
examples and exceptional exclusion are reproofs of assertions in
the supplied source. No rank upper bound for $q\ne3$ and no
solution of an open conjecture is asserted here.
\end{document}
